% =====================================================================
%  1. Introduction et problématique
% =====================================================================
\clearpage
\section{Problématique et traduction en tâche d'apprentissage}
\label{sec:intro}

\subsection{Contexte industriel}

Le soudage des aciers intervient dans des secteurs dont le poids économique se
chiffre en milliards d'euros, par exemple l'assemblage des mâts et pieux
d'éoliennes. Une soudure défaillante peut compromettre toute une structure ;
il faut donc garantir que le métal déposé possède la résistance et la ténacité
attendues. Aujourd'hui, la connaissance sur la qualité des soudures se transmet
surtout d'expert à expert, ce qui rend les industriels dépendants de quelques
spécialistes. L'énoncé du projet propose d'exploiter les données pour
\emph{extraire et homogénéiser} cette connaissance experte, et éventuellement
découvrir de nouveaux motifs.

\subsection{Qu'est-ce que la « qualité » d'une soudure ?}

La qualité n'est pas une grandeur mesurée directement dans la base. Elle se
manifeste à travers plusieurs \emph{propriétés mécaniques} mesurées par des
essais normalisés :
\begin{itemize}
  \item la \textbf{limite d'élasticité} (\YS) : contrainte à partir de laquelle le métal se déforme de façon permanente ;
  \item la \textbf{résistance maximale à la traction} (\UTS) : contrainte maximale supportée avant la rupture ;
  \item l'\textbf{allongement} : déformation atteinte à la rupture, qui renseigne sur la ductilité ;
  \item l'\textbf{énergie Charpy} : énergie absorbée lors de la rupture d'une éprouvette entaillée à une température donnée, qui renseigne sur la ténacité.
\end{itemize}
Ces propriétés décrivent des aspects différents et parfois antagonistes : un
métal très résistant est souvent moins ductile (\cref{subsec:eda_cibles}).
Fusionner ces mesures en un score unique \og bon/mauvais \fg{} demanderait un
cahier des charges (seuils par application, normes visées) que la base ne fournit
pas.

\begin{decision}
Nous ne créons \textbf{aucune classe arbitraire} \og bonne soudure / mauvaise soudure \fg{}.
Chaque propriété mesurée est prédite séparément, dans son unité physique. Ce choix
conserve l'information quantitative, évite d'inventer des seuils et permet à un
utilisateur d'appliquer ensuite ses propres critères d'acceptation.
\end{decision}

\subsection{Formulation comme problème d'apprentissage}

Pour chaque propriété cible $t \in \{\YS, \UTS, \text{allongement}, \text{Charpy}\}$,
on cherche une fonction $f_t$ qui associe aux \emph{conditions de fabrication}
d'une soudure $\bm{x}$ (composition chimique, paramètres électriques, apport
thermique, températures, traitement thermique, procédé) une estimation de la
propriété $y_t$ :
\begin{equation}
  \hat{y}_t = f_t(\bm{x}), \qquad \bm{x} \in \mathbb{R}^{p}\times\{\text{procédés}\},\quad y_t \in \mathbb{R}.
\end{equation}
Il s'agit donc de \textbf{régressions supervisées}. Pour l'énergie Charpy, la
température d'essai $T$ fait partie des entrées, car l'énergie dépend fortement
de la température : on prédit $E(T)$.

La base est \textbf{partiellement étiquetée} : une ligne peut renseigner la
composition et le procédé sans renseigner la propriété cible. Pour chaque cible,
on dispose ainsi d'un ensemble étiqueté $\mathcal{D}_L = \{(\bm{x}_i, y_i)\}$ et
d'un ensemble non étiqueté $\mathcal{D}_U = \{\bm{x}_j\}$. Cette situation motive
l'étude d'une méthode \textbf{semi-supervisée}, qui cherche à exploiter
$\mathcal{D}_U$ en plus de $\mathcal{D}_L$ (\cref{subsec:ssl_biblio}).

\begin{intuition}
Dans un usage réel, on voudrait estimer les propriétés d'une soudure \emph{avant}
de réaliser des essais mécaniques coûteux, à partir de sa recette (composition,
réglages, traitement). Les entrées autorisées sont donc uniquement celles
connues à ce moment-là : une autre propriété mécanique, la microstructure
observée après coup ou l'identifiant de la publication ne peuvent pas servir
d'entrée (\cref{subsec:entrees_exclues}).
\end{intuition}

\subsection{Démarche suivie}

Le projet est découpé en sept notebooks autonomes, exécutés dans l'ordre. La
\cref{fig:demarche} présente leur enchaînement et la section du rapport qui
correspond à chacun.

\begin{figure}[htbp]
\centering
\begin{tikzpicture}[
  etape/.style={rectangle, rounded corners=2pt, draw=csbordeaux, line width=0.7pt,
                fill=csgrisclair, align=center, text width=3.3cm, minimum height=1.35cm,
                font=\footnotesize},
  fleche/.style={-{Stealth[length=2.2mm]}, line width=0.8pt, color=csgris},
  node distance=0.55cm and 0.45cm]
  \node[etape] (n1) {\textbf{01 --- EDA}\\ structure, absences,\\ cibles, ACP\\ {\color{csbordeaux}\S\ref{sec:eda}}};
  \node[etape, right=of n1] (n2) {\textbf{02 --- Préparation}\\ nettoyage, groupes,\\ cohorte, plis\\ {\color{csbordeaux}\S\ref{sec:prep}}};
  \node[etape, right=of n2] (n3) {\textbf{03 --- Supervisé}\\ 5 familles, CV\\ groupée, 13/itérative\\ {\color{csbordeaux}\S\ref{subsec:bench13}}};
  \node[etape, right=of n3] (n4) {\textbf{04 --- Diagnostics}\\ erreurs par série,\\ sensibilités 11/13/16\\ {\color{csbordeaux}\S\ref{sec:diag}}};
  \node[etape, below=of n4] (n5) {\textbf{05 --- CV aléatoire}\\ 16/médiane,\\ exports\\ {\color{csbordeaux}\S\ref{subsec:res16}}};
  \node[etape, left=of n5] (n6) {\textbf{06 --- CV groupée}\\ 16/médiane,\\ exports\\ {\color{csbordeaux}\S\ref{subsec:res16}}};
  \node[etape, left=of n6] (n7) {\textbf{07 --- COREG}\\ semi-supervisé\\ inductif\\ {\color{csbordeaux}\S\ref{sec:ssl}}};
  \node[etape, left=of n7, fill=white, draw=csgris] (n8) {\textbf{Synthèse}\\ comparaison,\\ recommandations\\ {\color{csbordeaux}\S\ref{sec:discussion}}};
  \draw[fleche] (n1) -- (n2);
  \draw[fleche] (n2) -- (n3);
  \draw[fleche] (n3) -- (n4);
  \draw[fleche] (n4) -- (n5);
  \draw[fleche] (n5) -- (n6);
  \draw[fleche] (n6) -- (n7);
  \draw[fleche] (n7) -- (n8);
\end{tikzpicture}
\caption{Enchaînement des étapes du projet. Chaque boîte correspond à un notebook
du dossier \texttt{notebooks/} ; la référence indique la section du rapport qui
le décrit. Le diagnostic (04) a motivé le passage de la configuration initiale
13 entrées/imputation itérative à la configuration 16 entrées/imputation médiane.}
\label{fig:demarche}
\end{figure}

Trois principes ont guidé toutes les étapes :
\begin{enumerate}
  \item \textbf{Traçabilité} : les fichiers bruts ne sont jamais modifiés ; chaque transformation conserve la valeur d'origine, son statut et sa provenance.
  \item \textbf{Absence de fuite d'information} : toute transformation apprise (moyennes, médianes, encodage, imputation, choix d'hyperparamètres) est ajustée sur le seul ensemble d'entraînement du pli considéré.
  \item \textbf{Honnêteté des conclusions} : chaque résultat est rattaché à sa population et à son protocole ; les limites sont explicitées plutôt que masquées.
\end{enumerate}

\subsection{Organisation du dépôt}

\begin{table}[htbp]
\centering
\small
\caption{Organisation du dépôt du projet. Les données brutes, les données dérivées
et les modèles sont produits localement et ne sont pas versionnés.}
\label{tab:depot}
\begin{tabularx}{\textwidth}{@{}l L@{}}
\toprule
Élément & Rôle \\
\midrule
\fichier{sujet.md}                   & Énoncé original, conservé tel quel \\
\fichier{specs/spec.md}              & Décisions scientifiques et justifications (protocole actif) \\
\fichier{scripts/download\_data.py}  & Téléchargement autonome et vérification SHA-256 de l'archive Cambridge \\
\fichier{notebooks/01\ldots07}       & Code d'analyse, de préparation, d'entraînement et d'évaluation \\
\fichier{data/raw/}                  & Fichiers sources inchangés (\texttt{welddb.data}, \texttt{.info}, \texttt{.tex}) \\
\fichier{data/processed/}            & Exports dérivés : table maître, provenance, plis, scores, prédictions \\
\fichier{models/}                    & 188 pipelines complètes sérialisées avec leurs métadonnées \\
\fichier{reports/}                   & Rapports intermédiaires (01 à 16) et figures historiques \\
\fichier{tests/}                     & 37 tests ciblés hors ligne \\
\fichier{rapport\_latex/}            & Le présent rapport et le script de génération de ses figures \\
\bottomrule
\end{tabularx}
\end{table}

\begin{aretenir}
La qualité d'une soudure est représentée par quatre propriétés mécaniques
mesurées. Le problème devient quatre \textbf{régressions supervisées}
indépendantes, sur des données \textbf{partiellement étiquetées} --- ce qui
justifie une étude semi-supervisée complémentaire.
\end{aretenir}
